% =================================================================
%  CSE 815 -- Machine Learning
%  Module 3, Week 7, Part 2
%  Anomaly Detection: the Gaussian Distribution, Density Estimation,
%  Evaluation, and Feature Choice
%  (the epsilon frame is a step-by-step beamer overlay)
% =================================================================
\documentclass[aspectratio=169,10pt]{beamer}
\usepackage{cse815}

\title[CSE 815 -- Week 7, Part 2]{Anomaly Detection}
\subtitle{The Gaussian Distribution, Density Estimation, Evaluation, and Choosing Features}
\author{Rokan Uddin Faruqui\\ \small Professor, Department of CSE}
\institute{\small University of Chittagong \\ \texttt{rufaruqui@cu.ac.bd}}
\date[Week 7 / Part 2]{Week 7 $\cdot$ Part 2}

\begin{document}

\begin{frame}[plain,noframenumbering]
  \titlepage
\end{frame}

% ---------------------------------------------------------------
\begin{frame}{From groups to outliers}
  \begin{columns}[T,onlytextwidth]
    \begin{column}{0.5\textwidth}
      \begin{block}{Part 1 gave us}
        \footnotesize
        K-means: find groups in unlabelled data by alternating assignment and centroid
        updates, minimising the distortion $J$.
      \end{block}
      \footnotesize
      Clustering asks what the typical points look like. Anomaly detection asks the
      opposite: which points look \emph{unlike} the rest? It learns what normal is from
      unlabelled data, then flags whatever is improbable under that model.
    \end{column}
    \begin{column}{0.46\textwidth}
      \begin{keyidea}[Today]
        \footnotesize
        \begin{enumerate}
          \item Finding unusual events
          \item The Gaussian (normal) distribution
          \item The anomaly detection algorithm
          \item Developing and evaluating a detector
          \item Anomaly detection vs.\ supervised learning
          \item Choosing what features to use
        \end{enumerate}
      \end{keyidea}
    \end{column}
  \end{columns}
\end{frame}

% ===============================================================
\section{Finding Unusual Events}

\begin{frame}{Example: aircraft engines}
  \begin{columns}[T,onlytextwidth]
    \begin{column}{0.52\textwidth}
      \centering {\scriptsize 250 of 6000 engines that passed inspection (synthetic)}\par
      \begin{tikzpicture}
        \begin{axis}[width=7.0cm,height=4.9cm,xmin=380,xmax=620,ymin=1.8,ymax=6.2,
          xlabel={$x_1$: heat generated ($^\circ$C)},ylabel={$x_2$: vibration (mm/s)},
          tick label style={font=\scriptsize},label style={font=\scriptsize}]
          \addplot[only marks,mark=*,mark size=1pt,cublue!60] coordinates {(528.4,3.772)(509.5,3.686)(485.1,4.993)(474.0,4.142)(503.3,4.110)(514.0,4.218)(457.4,3.908)(493.9,3.909)(498.5,4.284)(457.8,4.709)(476.9,4.195)(536.4,4.242)(470.0,3.937)(468.1,4.541)(500.5,2.936)(544.8,4.672)(504.4,4.018)(521.1,4.382)(500.6,3.523)(490.2,4.044)(517.2,3.194)(495.0,4.542)(496.4,3.521)(594.1,4.166)(508.2,4.034)(483.2,3.610)(466.3,4.308)(528.3,3.986)(478.4,3.658)(527.3,5.004)(422.9,4.180)(473.8,4.050)(491.9,3.752)(468.9,4.429)(529.8,3.680)(474.5,4.020)(508.5,3.188)(514.8,3.633)(501.3,3.942)(502.1,3.299)(494.8,3.837)(499.3,4.248)(496.5,3.926)(493.8,4.584)(466.0,4.203)(539.1,4.173)(495.7,3.860)(537.8,4.476)(517.2,4.388)(526.9,4.122)(467.7,4.173)(483.9,3.011)(526.7,2.948)(506.2,4.551)(467.9,4.268)(513.2,3.650)(474.8,3.619)(532.2,4.055)(492.3,3.839)(457.2,3.994)(509.6,5.191)(489.0,3.872)(514.2,4.722)(507.8,4.591)(525.1,3.773)(489.7,4.257)(484.0,4.399)(473.4,2.920)(514.8,3.846)(541.3,4.467)(507.7,2.860)(482.3,3.513)(491.6,4.051)(474.0,4.160)(516.1,4.114)(505.6,3.007)(516.5,3.128)(514.7,3.858)(466.6,4.203)(510.1,3.698)(462.8,3.623)(457.6,3.319)(478.5,4.343)(503.2,3.406)(500.4,4.258)(465.5,4.475)(490.7,3.468)(480.3,3.971)(527.0,4.762)(522.4,3.865)(536.2,3.854)(465.2,4.115)(486.6,4.442)(507.4,4.055)(453.4,4.721)(447.1,3.562)(483.5,3.794)(469.9,3.229)(523.5,4.661)(513.4,4.585)(479.5,4.936)(493.8,4.193)(475.6,2.845)(506.9,4.392)(517.1,3.842)(487.9,4.466)(516.8,3.780)(510.8,4.337)(503.2,4.209)(523.7,4.347)(523.6,4.358)(524.2,3.480)(520.5,4.020)(470.7,3.530)(523.9,3.631)(489.7,4.210)(472.2,3.981)(551.9,3.962)(467.5,3.901)(493.4,4.581)(496.4,3.518)(456.1,3.708)(495.3,2.749)(477.2,5.048)(558.8,3.448)(495.6,4.019)(480.7,4.508)(476.9,4.722)(492.5,4.461)(529.7,3.306)(516.5,4.253)(478.8,3.789)(498.0,4.149)(483.2,4.595)(495.0,3.886)(490.5,3.448)(497.4,4.137)(503.5,3.892)(447.0,3.903)(516.8,3.594)(509.8,3.969)(496.3,2.734)(509.9,4.055)(528.5,3.708)(517.4,3.560)(465.2,4.308)(482.3,4.069)(533.8,4.242)(547.3,3.253)(465.8,3.274)(513.1,3.120)(511.5,5.008)(512.9,3.934)(452.6,3.838)(498.8,3.875)(564.0,4.533)(495.9,4.344)(553.1,4.103)(503.3,3.777)(503.1,4.300)(512.8,4.480)(526.2,4.215)(504.0,3.817)(573.2,3.479)(503.5,4.184)(533.0,3.486)(478.1,3.846)(520.6,3.827)(550.2,4.183)(515.8,4.065)(495.7,4.688)(533.4,4.037)(506.0,4.342)(510.5,3.914)(450.4,3.962)(520.8,3.813)(471.1,3.342)(515.8,3.711)(514.9,4.504)(493.8,4.051)(515.0,4.252)(543.7,3.706)(524.2,3.930)(530.3,3.755)(461.7,3.761)(487.1,2.970)(518.1,3.617)(522.4,4.190)(523.1,3.885)(507.8,4.793)(498.8,3.809)(477.8,4.299)(518.7,3.685)(504.0,4.025)(511.3,3.650)(513.6,3.999)(472.0,3.688)(501.5,4.150)(560.0,3.990)(498.1,3.802)(503.3,4.190)(505.2,3.789)(480.7,4.675)(534.2,5.154)(487.0,2.933)(458.4,3.697)(494.0,3.646)(447.4,3.076)(515.6,4.200)(505.9,3.122)(505.1,3.271)(477.9,3.307)(511.6,3.504)(545.9,3.422)(493.9,3.586)(479.8,3.639)(547.4,3.922)(498.9,4.104)(474.5,4.466)(514.9,3.367)(486.7,3.512)(497.7,3.637)(482.8,4.232)(507.7,4.213)(546.2,3.826)(522.8,3.379)(472.3,3.392)(475.3,4.088)(498.1,4.308)(483.0,3.802)(532.4,3.449)(506.0,3.023)(502.3,2.962)(498.4,3.974)(527.6,5.122)(472.3,3.052)(515.3,4.610)(500.4,5.095)(493.7,4.552)(498.7,3.688)(518.1,4.263)(509.3,3.544)(543.8,3.805)(467.4,4.204)(500.3,4.821)(454.9,3.920)(492.8,4.723)(478.0,3.726)(489.8,4.381)(483.4,4.323)};
          \addplot[only marks,mark=x,mark size=3pt,curose,very thick] coordinates {(590,5.6)};
          \node[font=\tiny,curose,anchor=west] at (axis cs:594,5.6) {new engine?};
        \end{axis}
      \end{tikzpicture}
    \end{column}
    \begin{column}{0.44\textwidth}
      \footnotesize
      \begin{block}{The task}
        Data: features $\vx^{(1)},\dots,\vx^{(m)}$ of engines already built --- almost all
        of them fine. A new engine $\vx_{\text{test}}$ comes off the line. Is it like the
        others, or should it be inspected?
      \end{block}
      \begin{defnbox}[Density estimation]
        \small Build a model $p(\vx)$ of how probable each feature vector is. Then
        \[ p(\vx_{\text{test}})<\varepsilon\ \Rightarrow\ \text{flag as an anomaly}. \]
      \end{defnbox}
      \begin{block}{Applications}
        Fraud detection (logins, purchases, typing speed); manufacturing (engines, circuit
        boards); monitoring computers in a data centre (memory, disk, CPU).
      \end{block}
    \end{column}
  \end{columns}
\end{frame}

% ===============================================================
\section{The Gaussian Distribution}

\begin{frame}{The Gaussian (normal) distribution}
  \begin{columns}[T,onlytextwidth]
    \begin{column}{0.5\textwidth}
      \centering
      \begin{tikzpicture}
        \begin{axis}[width=6.6cm,height=4.6cm,xmin=-4,xmax=6,ymin=0,ymax=0.85,
          xlabel={$x$},ylabel={$p(x)$},tick label style={font=\scriptsize},
          label style={font=\scriptsize},
          legend style={font=\tiny,at={(0.98,0.98)},anchor=north east},legend cell align=left]
          \addplot[cublue,very thick,domain=-4:6,samples=120]{exp(-x^2/2)/sqrt(2*pi)};
          \addlegendentry{$\mu=0,\ \sigma=1$}
          \addplot[curose,very thick,domain=-4:6,samples=120]{exp(-x^2/(2*0.25))/sqrt(2*pi*0.25)};
          \addlegendentry{$\mu=0,\ \sigma=0.5$}
          \addplot[cuteal,very thick,domain=-4:6,samples=120]{exp(-x^2/8)/sqrt(8*pi)};
          \addlegendentry{$\mu=0,\ \sigma=2$}
          \addplot[cuamber,very thick,dashed,domain=-4:6,samples=120]{exp(-(x-3)^2/(2*0.25))/sqrt(2*pi*0.25)};
          \addlegendentry{$\mu=3,\ \sigma=0.5$}
        \end{axis}
      \end{tikzpicture}
    \end{column}
    \begin{column}{0.46\textwidth}
      \begin{defnbox}[Gaussian density]
        \small
        \[ p(x;\mu,\sigma^2)=\frac{1}{\sqrt{2\pi}\,\sigma}\exp\!\Bigl(-\frac{(x-\mu)^2}{2\sigma^2}\Bigr) \]
        $\mu$: the centre. $\sigma$: the spread (standard deviation); $\sigma^2$ the
        variance.
      \end{defnbox}
      \footnotesize
      The area under each curve is $1$. Halve $\sigma$ and the bell becomes half as wide
      and twice as tall.
      \begin{pitfall}[A density is not a probability]
        \footnotesize With $\sigma=0.5$ the peak is $0.798$; with $\sigma=0.1$ it is
        $3.99$. Values above $1$ are fine --- only areas are probabilities.
      \end{pitfall}
    \end{column}
  \end{columns}
\end{frame}

\begin{frame}{Estimating $\mu$ and $\sigma^2$ from data}
  \begin{columns}[T,onlytextwidth]
    \begin{column}{0.5\textwidth}
      \begin{defnbox}[Parameter estimates]
        \small
        \[ \mu=\frac1m\sum_{i=1}^{m}x^{(i)},\qquad
           \sigma^2=\frac1m\sum_{i=1}^{m}\bigl(x^{(i)}-\mu\bigr)^2 \]
        These are the maximum-likelihood estimates. Statisticians often divide by $m-1$
        instead; for the thousands of examples typical here the difference is negligible.
      \end{defnbox}
    \end{column}
    \begin{column}{0.46\textwidth}
      \footnotesize
      \begin{block}{The 6000 engines}
        \renewcommand{\arraystretch}{1.2}
        \begin{tabular}{@{}lcc@{}}
          \toprule
          & $\mu_j$ & $\sigma_j$\\
          \midrule
          $x_1$, heat ($^\circ$C)     & $499.79$ & $25.12$\\
          $x_2$, vibration (mm/s)     & $3.993$ & $0.495$\\
          \bottomrule
        \end{tabular}
      \end{block}
      \begin{keyidea}
        Fitting a Gaussian is two averages. Nothing is trained by gradient descent; there
        are no labels.
      \end{keyidea}
    \end{column}
  \end{columns}
\end{frame}

% ===============================================================
\section{The Anomaly Detection Algorithm}

\begin{frame}{Density estimation with $n$ features}
  \begin{columns}[T,onlytextwidth]
    \begin{column}{0.52\textwidth}
      \begin{defnbox}[Anomaly detection algorithm]
        \small
        \begin{enumerate}
          \item Choose features $x_j$ that might indicate anomalies.
          \item Fit $\mu_j$, $\sigma_j^2$ for each feature $j=1,\dots,n$ on the training set.
          \item For a new $\vx$, compute
          \[ p(\vx)=\prod_{j=1}^{n}p(x_j;\mu_j,\sigma_j^2). \]
          \item Flag $\vx$ as an anomaly if $p(\vx)<\varepsilon$.
        \end{enumerate}
      \end{defnbox}
      \footnotesize
      Multiplying the densities treats the features as statistically independent. The
      algorithm works well in practice even when they are not quite.
    \end{column}
    \begin{column}{0.44\textwidth}
      \footnotesize
      \begin{block}{Three new engines}
        \renewcommand{\arraystretch}{1.2}
        \begin{tabular}{@{}lccc@{}}
          \toprule
          $\vx$ & $p(x_1)$ & $p(x_2)$ & $p(\vx)$\\
          \midrule
          $(505,4.1)$ & $0.0155$ & $0.787$ & $1.2\times10^{-2}$\\
          $(560,4.9)$ & $0.0009$ & $0.151$ & $1.4\times10^{-4}$\\
          $(590,5.6)$ & $0.000025$ & $0.0042$ & $1.0\times10^{-7}$\\
          \bottomrule
        \end{tabular}
      \end{block}
      With $\varepsilon=10^{-5}$: the first two are normal; the third, $3.6\sigma$ hot
      and $3.2\sigma$ shaky, is flagged.
      \begin{keyidea}
        $p(\vx)$ is small when \emph{any} feature is far from its mean --- one improbable
        reading is enough.
      \end{keyidea}
    \end{column}
  \end{columns}
\end{frame}

\begin{frame}{Turning the threshold $\varepsilon$}
  \begin{columns}[T,onlytextwidth]
    \begin{column}{0.52\textwidth}
      \centering
      \begin{tikzpicture}
        \begin{axis}[width=7.0cm,height=5.0cm,xmin=380,xmax=620,ymin=1.8,ymax=6.2,
          view={0}{90},xlabel={heat ($^\circ$C)},ylabel={vibration (mm/s)},
          tick label style={font=\scriptsize},label style={font=\scriptsize}]
          \addplot[only marks,mark=*,mark size=0.8pt,cublue!45,forget plot] coordinates {(528.4,3.772)(509.5,3.686)(485.1,4.993)(474.0,4.142)(503.3,4.110)(514.0,4.218)(457.4,3.908)(493.9,3.909)(498.5,4.284)(457.8,4.709)(476.9,4.195)(536.4,4.242)(470.0,3.937)(468.1,4.541)(500.5,2.936)(544.8,4.672)(504.4,4.018)(521.1,4.382)(500.6,3.523)(490.2,4.044)(517.2,3.194)(495.0,4.542)(496.4,3.521)(594.1,4.166)(508.2,4.034)(483.2,3.610)(466.3,4.308)(528.3,3.986)(478.4,3.658)(527.3,5.004)(422.9,4.180)(473.8,4.050)(491.9,3.752)(468.9,4.429)(529.8,3.680)(474.5,4.020)(508.5,3.188)(514.8,3.633)(501.3,3.942)(502.1,3.299)(494.8,3.837)(499.3,4.248)(496.5,3.926)(493.8,4.584)(466.0,4.203)(539.1,4.173)(495.7,3.860)(537.8,4.476)(517.2,4.388)(526.9,4.122)(467.7,4.173)(483.9,3.011)(526.7,2.948)(506.2,4.551)(467.9,4.268)(513.2,3.650)(474.8,3.619)(532.2,4.055)(492.3,3.839)(457.2,3.994)(509.6,5.191)(489.0,3.872)(514.2,4.722)(507.8,4.591)(525.1,3.773)(489.7,4.257)(484.0,4.399)(473.4,2.920)(514.8,3.846)(541.3,4.467)(507.7,2.860)(482.3,3.513)(491.6,4.051)(474.0,4.160)(516.1,4.114)(505.6,3.007)(516.5,3.128)(514.7,3.858)(466.6,4.203)(510.1,3.698)(462.8,3.623)(457.6,3.319)(478.5,4.343)(503.2,3.406)(500.4,4.258)(465.5,4.475)(490.7,3.468)(480.3,3.971)(527.0,4.762)(522.4,3.865)(536.2,3.854)(465.2,4.115)(486.6,4.442)(507.4,4.055)(453.4,4.721)(447.1,3.562)(483.5,3.794)(469.9,3.229)(523.5,4.661)(513.4,4.585)(479.5,4.936)(493.8,4.193)(475.6,2.845)(506.9,4.392)(517.1,3.842)(487.9,4.466)(516.8,3.780)(510.8,4.337)(503.2,4.209)(523.7,4.347)(523.6,4.358)(524.2,3.480)(520.5,4.020)(470.7,3.530)(523.9,3.631)(489.7,4.210)(472.2,3.981)(551.9,3.962)(467.5,3.901)(493.4,4.581)(496.4,3.518)(456.1,3.708)(495.3,2.749)(477.2,5.048)(558.8,3.448)(495.6,4.019)(480.7,4.508)(476.9,4.722)(492.5,4.461)(529.7,3.306)(516.5,4.253)(478.8,3.789)(498.0,4.149)(483.2,4.595)(495.0,3.886)(490.5,3.448)(497.4,4.137)(503.5,3.892)(447.0,3.903)(516.8,3.594)(509.8,3.969)(496.3,2.734)(509.9,4.055)(528.5,3.708)(517.4,3.560)(465.2,4.308)(482.3,4.069)(533.8,4.242)(547.3,3.253)(465.8,3.274)(513.1,3.120)(511.5,5.008)(512.9,3.934)(452.6,3.838)(498.8,3.875)(564.0,4.533)(495.9,4.344)(553.1,4.103)(503.3,3.777)(503.1,4.300)(512.8,4.480)(526.2,4.215)(504.0,3.817)(573.2,3.479)(503.5,4.184)(533.0,3.486)(478.1,3.846)(520.6,3.827)(550.2,4.183)(515.8,4.065)(495.7,4.688)(533.4,4.037)(506.0,4.342)(510.5,3.914)(450.4,3.962)(520.8,3.813)(471.1,3.342)(515.8,3.711)(514.9,4.504)(493.8,4.051)(515.0,4.252)(543.7,3.706)(524.2,3.930)(530.3,3.755)(461.7,3.761)(487.1,2.970)(518.1,3.617)(522.4,4.190)(523.1,3.885)(507.8,4.793)(498.8,3.809)(477.8,4.299)(518.7,3.685)(504.0,4.025)(511.3,3.650)(513.6,3.999)(472.0,3.688)(501.5,4.150)(560.0,3.990)(498.1,3.802)(503.3,4.190)(505.2,3.789)(480.7,4.675)(534.2,5.154)(487.0,2.933)(458.4,3.697)(494.0,3.646)(447.4,3.076)(515.6,4.200)(505.9,3.122)(505.1,3.271)(477.9,3.307)(511.6,3.504)(545.9,3.422)(493.9,3.586)(479.8,3.639)(547.4,3.922)(498.9,4.104)(474.5,4.466)(514.9,3.367)(486.7,3.512)(497.7,3.637)(482.8,4.232)(507.7,4.213)(546.2,3.826)(522.8,3.379)(472.3,3.392)(475.3,4.088)(498.1,4.308)(483.0,3.802)(532.4,3.449)(506.0,3.023)(502.3,2.962)(498.4,3.974)(527.6,5.122)(472.3,3.052)(515.3,4.610)(500.4,5.095)(493.7,4.552)(498.7,3.688)(518.1,4.263)(509.3,3.544)(543.8,3.805)(467.4,4.204)(500.3,4.821)(454.9,3.920)(492.8,4.723)(478.0,3.726)(489.8,4.381)(483.4,4.323)};
\only<1>{\costcontour{contour/draw color=cuteal,very thick,forget plot,domain=380:620,y domain=1.8:6.2}{5.0975}{((x-499.7865)/25.1190)^2+((y-3.99277)/0.49535)^2}
  \addplot[only marks,mark=o,mark size=1.6pt,cuamber,thick,forget plot] coordinates {(556.1,3.693)(573.9,3.996)(442.1,4.813)(557.9,3.786)(445.2,2.889)(461.2,2.915)(508.0,5.299)(432.1,3.592)(522.0,2.528)(564.5,3.870)(445.3,3.014)(499.2,5.128)(549.0,3.402)(509.1,5.108)(452.8,4.725)(458.2,4.763)(559.9,4.680)(500.5,5.119)(448.7,3.444)(522.7,2.719)(472.6,5.044)(533.1,4.939)(443.5,3.396)(515.7,5.503)(509.0,5.208)(444.6,3.568)(528.7,2.988)(562.7,3.890)(506.2,2.821)(547.8,2.539)(425.3,3.649)(528.0,5.299)(438.5,4.418)(430.7,4.791)(555.2,4.742)(469.3,5.011)(569.1,3.804)(474.7,2.913)(473.2,2.812)(446.0,4.402)(488.7,5.363)(518.1,5.392)(506.3,5.106)(554.4,4.321)(463.6,2.778)(433.2,4.012)(440.7,4.473)(499.6,2.570)(506.3,2.763)(572.7,4.381)(420.8,4.283)(513.4,5.109)(484.4,5.191)(500.8,2.715)(497.1,2.370)(445.8,3.142)(535.5,3.076)(449.0,3.430)(469.2,4.994)(565.6,4.185)(446.9,3.496)(437.4,4.309)(428.7,4.334)(534.2,5.024)(545.0,4.768)(545.6,4.724)(494.3,2.694)(559.2,3.800)(412.7,4.421)(487.2,2.785)(507.9,5.407)(491.1,2.853)(431.3,4.231)(550.9,2.996)(532.5,4.976)(456.9,4.929)(539.6,3.124)(544.3,5.241)(563.2,4.395)(504.3,2.658)(561.6,3.861)(556.4,4.777)(544.0,4.918)(464.0,5.251)(534.1,5.050)(549.1,4.804)(460.9,3.161)(452.5,2.878)(542.4,4.792)(478.2,2.476)(428.1,3.944)(470.0,4.952)(482.0,5.234)(548.0,2.455)(436.0,4.290)(429.9,4.039)(489.7,2.863)(520.8,2.863)(458.4,5.108)(429.7,5.119)(446.6,3.590)(550.0,2.520)(559.9,3.839)(559.5,3.439)(462.9,4.861)(470.6,2.898)(433.4,3.537)(504.0,5.221)(428.6,3.710)(423.9,2.816)(554.8,4.664)(443.6,3.681)(465.9,3.055)(557.9,3.345)(488.6,5.260)(438.9,4.625)(454.3,5.271)(546.6,3.090)(456.0,5.004)(442.3,4.627)(475.0,5.178)(442.0,4.451)(568.7,3.219)(490.5,5.172)(549.8,2.997)(566.0,4.247)(562.2,3.994)(564.1,4.637)(490.6,2.818)(474.1,2.705)(442.4,4.799)(542.0,3.106)(536.0,3.047)(562.4,3.910)(565.6,5.494)(548.9,4.565)(575.4,3.712)(535.6,5.067)(499.2,5.238)(426.9,3.837)(533.9,3.066)(507.0,2.794)(499.2,5.278)(508.6,5.596)(563.2,3.704)(439.7,3.917)(489.4,5.361)(477.2,5.166)(554.1,3.467)(534.3,2.775)(483.3,5.660)(457.6,4.804)(541.0,2.981)(489.6,2.778)(549.2,2.449)(444.6,4.740)};}%
\only<2>{\costcontour{contour/draw color=cuteal,very thick,forget plot,domain=380:620,y domain=1.8:6.2}{9.7027}{((x-499.7865)/25.1190)^2+((y-3.99277)/0.49535)^2}
  \addplot[only marks,mark=o,mark size=1.6pt,cuamber,thick,forget plot] coordinates {(547.8,2.539)(430.7,4.791)(420.8,4.283)(497.1,2.370)(412.7,4.421)(478.2,2.476)(548.0,2.455)(429.7,5.119)(550.0,2.520)(423.9,2.816)(454.3,5.271)(568.7,3.219)(565.6,5.494)(508.6,5.596)(483.3,5.660)(549.2,2.449)};}%
\only<3>{\costcontour{contour/draw color=cuteal,very thick,forget plot,domain=380:620,y domain=1.8:6.2}{12.1106}{((x-499.7865)/25.1190)^2+((y-3.99277)/0.49535)^2}
  \addplot[only marks,mark=o,mark size=1.6pt,cuamber,thick,forget plot] coordinates {(547.8,2.539)(412.7,4.421)(548.0,2.455)(429.7,5.119)(550.0,2.520)(423.9,2.816)(565.6,5.494)(549.2,2.449)};}%
\only<4>{\costcontour{contour/draw color=cuteal,very thick,forget plot,domain=380:620,y domain=1.8:6.2}{14.3078}{((x-499.7865)/25.1190)^2+((y-3.99277)/0.49535)^2}
  \addplot[only marks,mark=o,mark size=1.6pt,cuamber,thick,forget plot] coordinates {(423.9,2.816)(565.6,5.494)};}%
\only<5>{\costcontour{contour/draw color=cuteal,very thick,forget plot,domain=380:620,y domain=1.8:6.2}{16.7158}{((x-499.7865)/25.1190)^2+((y-3.99277)/0.49535)^2}}%
          \addplot[only marks,mark=x,mark size=2.6pt,curose,very thick,forget plot] coordinates {(386.4,3.831)(599.2,3.362)(558.5,5.493)(609.5,3.268)(464.5,5.733)(411.0,4.930)(540.4,2.480)(417.4,5.060)(415.1,3.453)(590.4,4.316)};
        \end{axis}
      \end{tikzpicture}\par
      {\scriptsize \textcolor{cuteal}{teal}: the boundary $p(\vx)=\varepsilon$\\
      \textcolor{curose}{$\times$}: the 10 faulty validation engines \quad
      \textcolor{cuamber}{$\circ$}: good engines flagged}
    \end{column}
    \begin{column}{0.44\textwidth}
      \footnotesize
      \begin{block}{Validation set: 2000 good, 10 faulty}
        \renewcommand{\arraystretch}{1.1}\setlength{\tabcolsep}{3pt}
        \begin{tabular}{@{}lcccccc@{}}
          \toprule
          $\varepsilon$ & TP & FP & FN & $P$ & $R$ & $F_1$\\
          \midrule
          \alert<1>{$10^{-3}$} & \alert<1>{10} & \alert<1>{156} & \alert<1>{0} & \alert<1>{$0.060$} & \alert<1>{$1.0$} & \alert<1>{$0.114$}\\
          \alert<2>{$10^{-4}$} & \alert<2>{10} & \alert<2>{16} & \alert<2>{0} & \alert<2>{$0.385$} & \alert<2>{$1.0$} & \alert<2>{$0.556$}\\
          \alert<3>{$3\times10^{-5}$} & \alert<3>{9} & \alert<3>{8} & \alert<3>{1} & \alert<3>{$0.529$} & \alert<3>{$0.9$} & \alert<3>{$0.667$}\\
          \alert<4>{$10^{-5}$} & \alert<4>{7} & \alert<4>{2} & \alert<4>{3} & \alert<4>{$0.778$} & \alert<4>{$0.7$} & \alert<4>{$0.737$}\\
          \alert<5>{$3\times10^{-6}$} & \alert<5>{3} & \alert<5>{0} & \alert<5>{7} & \alert<5>{$1.000$} & \alert<5>{$0.3$} & \alert<5>{$0.462$}\\
          \bottomrule
        \end{tabular}
      \end{block}
      \begin{minipage}[t][7ex][t]{\linewidth}
\only<1>{$\varepsilon=10^{-3}$: flags 10 of 10 faulty and 156 of 2000 good validation engines. $F_1=0.114$.}%
\only<2>{$\varepsilon=10^{-4}$: flags 10 of 10 faulty and 16 of 2000 good validation engines. $F_1=0.556$.}%
\only<3>{$\varepsilon=3\times10^{-5}$: flags 9 of 10 faulty and 8 of 2000 good validation engines. $F_1=0.667$.}%
\only<4>{$\varepsilon=10^{-5}$: flags 7 of 10 faulty and 2 of 2000 good validation engines. $F_1=0.737$.}%
\only<5>{$\varepsilon=3\times10^{-6}$: flags 3 of 10 faulty and 0 of 2000 good validation engines. $F_1=0.462$.}%
      \end{minipage}
      \begin{keyidea}
        The boundary of $p(\vx)=\varepsilon$ is an ellipse (the real level set, drawn
        from the formula). Smaller $\varepsilon$ widens it: fewer false alarms, more
        missed faults.
      \end{keyidea}
    \end{column}
  \end{columns}
\end{frame}

% ===============================================================
\section{Developing and Evaluating a Detector}

\begin{frame}{Real-number evaluation, with a few labels}
  \begin{columns}[T,onlytextwidth]
    \begin{column}{0.5\textwidth}
      \begin{block}{The engines: 10\,000 good, 20 faulty}
        \footnotesize
        \renewcommand{\arraystretch}{1.2}
        \begin{tabular}{@{}lll@{}}
          \toprule
          set & good & faulty\\
          \midrule
          training      & $6000$ & $0$ (assumed)\\
          cross-validation & $2000$ & $10$\\
          test          & $2000$ & $10$\\
          \bottomrule
        \end{tabular}\par
        \vspace{1mm}
        Train on unlabelled, mostly-normal data. Use the few labelled anomalies only to
        choose $\varepsilon$ (and features), then to report.
      \end{block}
      \footnotesize
      A few anomalies slipping into the training set is harmless: they barely move $\mu$
      and $\sigma$.
    \end{column}
    \begin{column}{0.46\textwidth}
      \footnotesize
      \begin{block}{Procedure}
        \begin{enumerate}
          \item Fit $p(\vx)$ on the training set.
          \item On validation, predict $y=1$ if $p(\vx)<\varepsilon$.
          \item Classes are skewed ($10$ in $2010$): measure $P$, $R$, $F_1$ (Week 6), not
                accuracy.
          \item Keep the best-$F_1$ $\varepsilon$: $10^{-5}$, $F_1=0.737$.
        \end{enumerate}
      \end{block}
      \begin{keyidea}[Test set, once]
        At $\varepsilon=10^{-5}$: TP $=8$, FP $=0$, FN $=2$; precision
        $1.0$, recall $0.8$, $F_1=0.889$.
      \end{keyidea}
      \begin{pitfall}[Very few anomalies]
        \footnotesize With only $2$ faults, use both for validation and skip the test set
        --- the result is then optimistic.
      \end{pitfall}
    \end{column}
  \end{columns}
\end{frame}

% ===============================================================
\section{Anomaly Detection vs.\ Supervised Learning}

\begin{frame}{If we have some labels, why not train a classifier?}
  \footnotesize
  \renewcommand{\arraystretch}{1.35}
  \begin{tabular}{@{}L{6.6cm}L{6.8cm}@{}}
    \toprule
    \textbf{Anomaly detection} & \textbf{Supervised learning}\\
    \midrule
    Very few positives ($0$--$20$ is common), many negatives.
      & Enough positives and negatives to learn both.\\
    Many \emph{different kinds} of anomaly; future ones may look like none seen so far.
      & Future positives are likely to resemble the training positives.\\
    Fraud, where attackers invent new tricks; a never-before-seen engine fault; a
    misbehaving server.
      & Spam (the types recur); known manufacturing defects such as scratches; weather;
      diagnosing a known disease.\\
    \bottomrule
  \end{tabular}
  \vspace{3mm}
  \begin{columns}[T,onlytextwidth]
    \begin{column}{0.5\textwidth}
      \begin{keyidea}[The difference in one line]
        A classifier learns what positives look like. An anomaly detector learns what
        \emph{normal} looks like, and flags anything else.
      \end{keyidea}
    \end{column}
    \begin{column}{0.46\textwidth}
      \begin{block}{In practice}
        Many systems use both: a classifier for the known fraud patterns, an anomaly
        detector for the unknown ones.
      \end{block}
    \end{column}
  \end{columns}
\end{frame}

% ===============================================================
\section{Choosing What Features to Use}

\begin{frame}{Make features look Gaussian}
  \begin{columns}[T,onlytextwidth]
    \begin{column}{0.56\textwidth}
      \centering
      \begin{tikzpicture}
        \begin{axis}[name=h1,ybar,bar width=7pt,width=4.8cm,height=4.0cm,xmin=-0.25,xmax=6.25,
          ymin=0,title={\scriptsize raw $x$: skewed},tick label style={font=\tiny},
          ytick=\empty]
          \addplot[fill=cublue!60,draw=cublue] coordinates {(0.25,176)(0.75,310)(1.25,181)(1.75,139)(2.25,60)(2.75,41)(3.25,29)(3.75,22)(4.25,15)(4.75,4)(5.25,6)(5.75,3)};
        \end{axis}
        \begin{axis}[at={(h1.east)},anchor=west,xshift=4mm,ybar,bar width=7pt,width=4.8cm,
          height=4.0cm,xmin=-3.25,xmax=3.25,ymin=0,title={\scriptsize $\log x$: bell-shaped},
          tick label style={font=\tiny},ytick=\empty]
          \addplot[fill=cuteal!60,draw=cuteal] coordinates {(-2.75,1)(-2.25,8)(-1.75,21)(-1.25,63)(-0.75,159)(-0.25,234)(0.25,227)(0.75,171)(1.25,89)(1.75,19)(2.25,6)(2.75,2)};
        \end{axis}
      \end{tikzpicture}\par
      {\scriptsize 1000 synthetic readings; the raw histogram omits 14 values above 6.}
    \end{column}
    \begin{column}{0.4\textwidth}
      \footnotesize
      \begin{block}{Why it matters}
        The algorithm models each feature as a Gaussian. A long right tail makes large but
        perfectly normal values look improbable, and squeezes the real outliers together.
      \end{block}
      \begin{block}{Transformations to try}
        $\log x$, \ $\log(x+c)$, \ $\sqrt{x}$, \ $x^{1/3}$. Plot the histogram, choose by
        eye, and apply the \emph{same} transformation to the validation and test sets.
      \end{block}
    \end{column}
  \end{columns}
\end{frame}

\begin{frame}{Error analysis: invent the feature that exposes the miss}
  \begin{columns}[T,onlytextwidth]
    \begin{column}{0.5\textwidth}
      \centering {\scriptsize Data-centre computers (synthetic): traffic tracks CPU load}\par
      \begin{tikzpicture}
        \begin{axis}[width=6.4cm,height=4.6cm,xmin=10,xmax=100,ymin=0,ymax=1100,
          xlabel={CPU load (\%)},ylabel={network traffic},tick label style={font=\scriptsize},
          label style={font=\scriptsize},scaled y ticks=false]
          \addplot[only marks,mark=*,mark size=1pt,cublue!60] coordinates {(76.4,673.5)(23.8,276.0)(23.4,248.2)(50.4,500.0)(47.5,483.4)(58.9,541.3)(67.5,686.7)(81.1,703.3)(55.2,517.6)(76.4,780.2)(51.4,512.9)(76.0,824.9)(21.0,229.4)(45.8,389.5)(58.9,569.8)(67.9,768.1)(27.5,301.9)(75.9,764.5)(46.6,489.0)(89.4,986.1)(69.9,663.4)(53.6,534.2)(20.1,204.8)(84.7,887.7)(77.6,708.4)(80.1,864.7)(38.1,384.1)(80.5,821.2)(20.2,189.9)(82.3,748.7)(80.1,857.0)(61.6,624.6)(68.8,706.0)(35.0,328.6)(62.1,721.5)(49.2,509.6)(46.3,436.5)(36.1,341.1)(57.4,483.3)(23.0,243.9)(21.4,186.8)(71.1,810.7)(63.1,667.5)(33.4,379.6)(40.6,338.8)(21.2,197.2)(68.3,663.4)(67.5,643.9)(39.2,373.5)(71.3,662.3)(26.5,222.0)(20.0,210.7)(62.1,625.4)(62.8,583.1)(47.7,494.0)(55.5,508.3)(53.8,506.7)(81.6,824.3)(24.2,246.9)(57.9,534.4)(84.3,835.9)(27.9,288.3)(24.7,252.4)(53.6,521.7)(25.0,252.1)(71.9,703.5)(49.8,570.9)(43.5,423.7)(86.3,723.6)(59.7,600.5)(82.1,832.6)(74.8,770.7)(29.1,269.3)(77.2,723.2)(90.0,1024.8)(66.0,605.7)(24.9,263.0)(37.1,371.6)(24.3,235.9)(85.2,916.5)(34.3,379.6)(60.1,619.0)(61.7,606.8)(54.3,555.3)(56.1,578.5)(37.8,367.6)(68.9,738.5)(43.3,423.6)(68.5,613.7)(25.4,263.3)(37.9,371.6)(22.7,234.3)(75.4,659.8)(24.6,293.6)(71.2,658.1)(36.7,349.2)(34.8,347.8)(21.5,209.0)(74.5,663.3)(65.5,674.4)(28.7,292.8)(32.0,278.9)(83.7,953.0)(49.7,479.8)(35.9,347.1)(85.8,802.4)(75.1,734.1)(22.6,200.2)(59.6,628.1)(20.8,189.9)(39.7,392.2)(36.4,387.1)(79.0,690.9)(56.4,570.3)(71.0,690.6)(20.4,218.7)(77.9,843.9)(64.1,555.0)(24.4,231.7)(82.5,917.4)(57.7,617.6)(80.7,875.5)(86.0,909.2)(43.7,406.2)(36.8,373.7)};
          \addplot[only marks,mark=x,mark size=3pt,curose,very thick] coordinates {(85,250)};
          \node[font=\tiny,curose,anchor=east] at (axis cs:82,250) {stuck process?};
        \end{axis}
      \end{tikzpicture}
    \end{column}
    \begin{column}{0.46\textwidth}
      \footnotesize
      \begin{block}{The miss}
        High CPU with little traffic: a process stuck in a loop. Each feature alone is
        unremarkable ($1.5\sigma$ and $-1.4\sigma$), so $p(\vx)=4.5\times10^{-6}$ --- lower
        than for $79$ of the $500$ normal machines. No $\varepsilon$ flags it without
        flagging those $79$.
      \end{block}
      \begin{keyidea}[Add a feature]
        $x_3=\dfrac{\text{CPU load}}{\text{network traffic}}$. Normal machines:
        $0.101\pm0.009$. The stuck one: $0.34$, which is $28\sigma$ out. Now
        $p(\vx)\approx10^{-174}$, against at least $1.1\times10^{-7}$ for every normal
        machine.
      \end{keyidea}
      Error analysis for anomaly detection: find the anomalies $p(\vx)$ missed, ask what
      makes them unusual, and turn that into a feature.
    \end{column}
  \end{columns}
\end{frame}

% ===============================================================
\section{Wrap-up and Assessment}

\begin{frame}{Summary}
  \begin{columns}[T,onlytextwidth]
    \begin{column}{0.5\textwidth}
      \begin{block}{What to remember}
        \footnotesize
        \begin{itemize}
          \item Model normal data with $p(\vx)$; flag $p(\vx)<\varepsilon$.
          \item Fit one Gaussian per feature: $\mu_j$ and $\sigma_j^2$ are averages.
          \item Train on unlabelled data; use a few labelled anomalies to choose
                $\varepsilon$ by $F_1$, and to test.
          \item Prefer anomaly detection when positives are rare and varied; supervised
                learning when they are plentiful and alike.
          \item Transform skewed features towards a bell shape; build new features from
                the anomalies you miss.
        \end{itemize}
      \end{block}
    \end{column}
    \begin{column}{0.46\textwidth}
      \begin{block}{Formulas}
        \[ p(x;\mu,\sigma^2)=\tfrac{1}{\sqrt{2\pi}\sigma}e^{-(x-\mu)^2/2\sigma^2} \]
        \[ p(\vx)=\prod_{j=1}^{n}p(x_j;\mu_j,\sigma_j^2) \]
      \end{block}
      \begin{block}{Next}
        \footnotesize Week 8: recommender systems --- collaborative filtering and
        content-based filtering.
      \end{block}
    \end{column}
  \end{columns}
\end{frame}

\begin{frame}{Descriptive questions}
  \footnotesize
  \begin{enumerate}
    \item Explain anomaly detection by density estimation, and give three applications.
          \hfill\clo{4}
    \item State the Gaussian density. Sketch it for $(\mu,\sigma)=(0,1)$, $(0,0.5)$ and
          $(3,1)$, and explain why its value can exceed $1$. \hfill\clo{6}
    \item Give the estimates of $\mu$ and $\sigma^2$ from $m$ examples, and compute them for
          $x=\{2,4,4,4,5,5,7,9\}$. \hfill\clo{5}
    \item State the anomaly detection algorithm for $n$ features. What assumption does the
          product make? \hfill\clo{5}
    \item Using the lecture's engine parameters, compute $p(\vx)$ for $\vx=(450,3.0)$ and
          decide whether it is an anomaly at $\varepsilon=10^{-5}$. \hfill\clo{8}
    \item Explain how the data are split to develop an anomaly detector, and why accuracy
          is a poor way to choose $\varepsilon$. \hfill\clo{6}
  \end{enumerate}
\end{frame}

\begin{frame}{Descriptive questions (continued)}
  \footnotesize
  \begin{enumerate}\setcounter{enumi}{6}
    \item From the validation table, explain why $F_1$ is best at $\varepsilon=10^{-5}$
          rather than at the largest or smallest $\varepsilon$. \hfill\clo{5}
    \item Compare anomaly detection with supervised learning, and classify: (a) detecting
          new kinds of credit-card fraud; (b) filtering spam; (c) finding dented car
          panels in a factory with $5000$ labelled dents. \hfill\clo{6}
    \item Why should features be roughly Gaussian? Name three transformations, and explain
          why the same one must be applied to every data split. \hfill\clo{5}
    \item Explain error analysis for anomaly detection using the stuck-process example.
          \hfill\clo{6}
    \item Show that the level set $p(\vx)=\varepsilon$ for two features is an ellipse, and
          find its semi-axes for the engines at $\varepsilon=10^{-5}$. \hfill\clo{8}
    \item Why do a few anomalies in the training set do little harm? \hfill\clo{4}
  \end{enumerate}
\end{frame}

\begin{frame}{Design questions}
  \footnotesize
  \begin{enumerate}
    \item \textbf{Design a fraud detector.} A mobile financial service in Bangladesh wants
          to flag unusual transactions. Design the features (including at least one ratio),
          the training data, how you will obtain labelled frauds for validation, and how
          you choose $\varepsilon$ given the cost of a blocked genuine payment.
          \hfill\clo{12}
    \item \textbf{Design for a garment factory.} Design an anomaly detector for sewing
          machines using vibration and power sensors, including how you would detect a
          machine that drifts slowly over weeks. \hfill\clo{10}
    \item \textbf{Design the decision.} You have $40$ labelled defects and $50\,000$ good
          circuit boards. Design an experiment to decide between anomaly detection and a
          supervised classifier. \hfill\clo{12}
    \item \textbf{Design a feature-transform check.} Design an automatic procedure that
          picks a transformation for each feature before fitting $p(\vx)$, and a test
          that it improves validation $F_1$. \hfill\clo{10}
  \end{enumerate}
\end{frame}

\begin{frame}{Reading}
  \begin{block}{Primary text}
    G.~James, D.~Witten, T.~Hastie and R.~Tibshirani,
    \emph{An Introduction to Statistical Learning}, 2nd ed., Springer, 2021.\\[2pt]
    \hfill{\small$\rightarrow$ \textbf{\S4.4.1--4.4.2} --- Gaussian densities and their
    parameter estimates; \textbf{\S4.4.2} --- thresholds, sensitivity and specificity.}
  \end{block}
  \begin{block}{Supplementary}
    C.~M.~Bishop, \emph{Pattern Recognition and Machine Learning}, Springer, 2006.\\[2pt]
    \hfill{\small$\rightarrow$ \textbf{\S1.2.4} --- the Gaussian distribution and its
    maximum-likelihood estimates; \textbf{\S2.3} --- the multivariate Gaussian, whose
    level sets are ellipses.}
  \end{block}
  \vspace{1mm}
  \begin{center}
    \small\color{cugray}
    Lab (CSE 816, Week 7): anomaly detection.
  \end{center}
\end{frame}

\end{document}
